\section{Объект и методы исследования}\label{sec:methods}

\subsection{Вычислительная схема: ансамблевая векторизация}

Прямолинейная реализация сканирования требует обучить $n^2$ сетей независимо;
при $n=512$ и $T=400$ это $10^8$ итераций и, в наивном цикле, десятки часов
машинного времени. В настоящей работе используется другая схема: все узлы
сетки обучаются \emph{одновременно} как ансамбль независимых сетей.

Веса хранятся в тензорах формы $(G, h, d_{in})$ и $(G, d_{out}, h)$, где
$G$ --- число конфигураций в текущем блоке, а гиперпараметры --- в векторах
формы $(G,)$, транслируемых по осям весов. Прямой и обратный проход
записываются через пакетные свёртки индексов:
\begin{center}
\ttfamily\small
H = einsum('nd,ghd->gnh', X, W0);\quad
gW0 = einsum('gnh,nd->ghd', gh, X)
\end{center}
и так далее. Это сводит всю итерацию к нескольким вызовам BLAS и даёт
ускорение примерно на два порядка по сравнению с поэлементным циклом.
Расчёт ведётся блоками по $4096$ конфигураций, что ограничивает пик потребления
памяти примерно $200$~МБ.

Градиенты вычисляются вручную (без автоматического дифференцирования), что
обеспечивает полный контроль над порядком операций и, следовательно, побитовую
воспроизводимость. Конфигурации, достигшие переполнения, не исключаются из
расчёта, но их градиенты обнуляются через \texttt{nan\_to\_num}, чтобы
переполнение в одной конфигурации не портило соседние через общие
арифметические операции; факт переполнения при этом запоминается и определяет
режим NaN/Inf.

\subsection{Протокол экспериментов}

Общие параметры, если не указано иное: $d_{in}=d_{out}=h=16$; инициализация
весов нормальным распределением с масштабом $1.0$, одна и та же для всех узлов
сетки; число итераций $T=400$; тип данных \texttt{float64}; разрешение сетки
$256\times256$, для двух опорных карт --- $512\times512$.

Базовый набор данных --- одна обучающая пара ($|\mathcal{D}|=1$) со случайными
$x$ и $y$ из стандартного нормального распределения. Этот выбор унаследован от
работы~\cite{sohl2024}, где случай $|\mathcal{D}|=1$ рассматривается как
минимальный, в котором явление уже наблюдается, и не осложнён усреднением по
выборке. Для экспериментов с мини-батчами и со структурированными данными
используются наборы из 32 и 16 примеров соответственно.

Набор \texttt{sklearn digits} преобразуется следующим образом: изображения
$8\times8$ разворачиваются в векторы длины 64, центрируются, проецируются на
первые 16 главных компонент всей выборки и нормируются на единичное
стандартное отклонение; целевые векторы --- one-hot представление класса в
$\{-1,+1\}^{16}$. Такое преобразование позволяет использовать ту же
архитектуру, что и для случайных данных, и делает сравнение корректным.

\subsection{Выбор окон сканирования}

Порог устойчивости зависит от кривизны функции потерь, а та --- от активации,
объёма данных и оптимизатора. Поэтому единое окно по $\eta$ для всех
конфигураций непригодно: например, для ReLU при $|\mathcal{D}|=1$ во всём
диапазоне $\eta\in[10^{-2},10]$ доля сходимости составляет $0.99$, то есть
граница просто не попадает в окно. Окна подбирались так, чтобы граница
находилась внутри области сканирования, а доли режимов были нетривиальны.
Конкретные диапазоны указаны при каждом эксперименте; следствия этого выбора
для сопоставимости оценок обсуждаются в разделе~\ref{sec:limitations}.

\subsection{Последовательность зумов}

Центр зума выбирается автоматически, а не вручную: по опорной карте
$(\eta,\beta)$ вычисляется бинарная граница, к ней применяется усредняющий
фильтр $41\times41$, и центром объявляется точка максимума плотности границы,
удалённая не менее чем на 64 пикселя от краёв. Для базовой карты это дало
точку $\eta^\ast \approx 0.369$, $\beta^\ast \approx 0.300$. На уровне $k$
окно сжимается вдвое по $\log_{10}\eta$ и по $\beta$ относительно центра.

\subsection{Программная среда и ресурсы}

Реализация выполнена на Python 3.12 с использованием NumPy (линейная алгебра),
SciPy (морфологические операции), scikit-image (скелетонизация),
scikit-learn (набор digits) и Matplotlib (визуализация). PyTorch не
использовался: в целевой вычислительной среде он недоступен, а ансамблевая
векторизация на NumPy оказалась достаточной по производительности.
Вычисления проводились на одном ядре CPU; суммарное машинное время всех
экспериментов, вошедших в отчёт, составило \TotalHours~ч.

Воспроизводимость обеспечивается фиксацией: seed генератора случайных чисел
($42$ для данных, $42+1000$ для инициализации весов, $42+7$ для порядка
примеров в мини-батчах), типа точности, разрешения сетки, числа итераций и
порогов классификации режимов. Все параметры каждого запуска сохраняются в
JSON рядом с результатом, а числа в таблицах и в тексте настоящего отчёта
подставляются автоматически из этих файлов, что исключает расхождение между
текстом и данными.
